\scanpage{64}
(i) Si $\deleted{\unsure{\mathcal O}}\unsure{\underline{\mathbb W}|\underline W}^a(\unsure{\mathfrak X|\mathcal X})$ et $\underline{\mathbb W}^a(\unsure{\mathfrak Y|\mathcal Y})$ sont définis, il en est de $\unsure{=|\simeq}$ de $\mathbb W^a(\mathfrak X\times_{\mathbb W(S)}\mathfrak Y)$, et on a
\[
 \mathbb W^a(\mathfrak X\times_{\mathbb W(S)}\mathfrak Y)\simeq \mathbb W^a(\mathfrak X)\times_S \mathbb W^a(\mathfrak Y)
\]
et on \deleted{\unsure{a}} plus généralement \deleted{\unsure{fon|pour}} la formule analogue pour les limites projectives finies [\unsure{donc|d'où} $(\alpha)$ l'objet unité $(\beta)$ les produits fibrés, d'où comme conséquence $(\gamma)$ Noyaux]

\[
 W^a(\mathfrak X)\times_S T\simeq W_T^a(\mathfrak X\times_{\mathbb W(S)}\mathbb W(T))\tag*{(ii)}
\]
[l'existence du premier \unsure{mb|m\textsuperscript{b}} impliquant celle du second\,; la définition de $W_T$ est claire \unsure{\dots}]

De (i) résulte que le foncteur $\mathbb W^a$ transforme un préschéma en groupes, anneaux etc en préschéma du $=$ type.
% archive sheet number: 63
